\documentclass{article}
\usepackage[UTF8]{ctex}
\usepackage{amsmath,amssymb,amsfonts}
\usepackage{enumitem}
\usepackage{geometry}
\geometry{a4paper, margin=1in}
% === 错题本排版优化（不参与题目标识） ===
\setlength{\parindent}{0pt}
\setlength{\parskip}{0.35em plus 0.10em minus 0.08em}
\setlength{\emergencystretch}{3em}
\setlist{itemsep=0.25em, topsep=0.40em, parsep=0pt}
\allowdisplaybreaks[2]
% === 错题本排版优化结束 ===
\begin{document}
\section*{南京航空航天大学2026年高等代数真题}
\section*{93．南京航空航天大学2026年高等代数试卷（回忆版）}
（注．题目来自考生回忆，部分表述或数据可能有误）\\

\subsection{题目1}
1．只记得答案为 \(k(1,1,1,1)^{\mathrm{T}}, A X=k(1,1,1,1)^{\mathrm{T}}\) 的解空间。

\subsection{题目2}
2．\(V_{1} \neq V_{2}\) 时 \(a, b\) 的取值 \((a=2, b \in \mathbb{R}\) 或 \(b=1, a \in \mathbb{R}), \operatorname{dim}\left(V_{1} \cap V_{2}\right)=3\) ．

\subsection{题目3}
3．（1）基下矩阵，（2）\(A^{2025}\)（超难算，麻烦）（注：前三题类似往年真题，但计算量比往年大）

\subsection{题目4}
4．设矩阵 \(J=\left(\begin{array}{lll}0 & 1 & 0 \\ 0 & 0 & 1 \\ 1 & 0 & 0\end{array}\right), A=3 E+2 J+2 J^{2}\) ．

\begin{itemize}
  \item[（1）] 求 \(A\) 的最小多项式．
  \item[（2）] 求与 \(A\) 可交换的矩阵空间的基与维数．
\end{itemize}

\subsection{题目5}
5．解答如下问题：

\begin{itemize}
  \item[（1）] \(n\) 阶矩阵 \(A\) 正定，证明：\(f\left(x_{1}, x_{2}, \cdots, x_{n}\right)=\left|\begin{array}{cc}A & X \\ X^{\mathrm{T}} & 0\end{array}\right|\) 负定．
  \item[（2）] \(\alpha, \beta\) 为 \(n\) 维列向量，证明：\(\left|A+\alpha^{\mathrm{T}} \beta\right|=|A|+\alpha^{\mathrm{T}} A^{*} \beta\) 。
\end{itemize}

\subsection{题目6}
6．解答如下问题：

\begin{itemize}
  \item[（1）] 已知 \(\bar{A}^{\mathrm{T}} A=E\) ，证明：\(A\) 特征值模长为 1 。
  \item[（2）] 已知 \(A, B\) 均为 \(n\) 阶实对称矩阵，其中 \(n \geq 2\) ，且 \(A\) 正定，\(A B\) 的特征值全为 1 。证明：存在次数小于 \(n\) 的多项式 \(f(x)\) ，使得 \(B=f(A)\) 。
\end{itemize}

\subsection{题目7}
7．已知线性变换 \(\mathcal{T}(X)=A X B\) ，其中 \(A\) 为 \(m\) 阶方阵，\(B\) 为 \(n\) 阶方阵，\(A\) 的特征值为 \(\lambda_{i}, B\) 的特征值为 \(\mu_{j}\) 、证明：

\begin{itemize}
  \item[（1）] \(\lambda_{i} \mu_{j}\) 为 \(\mathcal{T}\) 的特征值。
  \item[（2）] 存在正整数 \(k\) 使得 \(\mathcal{T}^{k}=\mathcal{O}\) 等价于存在正整数 \(s\) 使得 \(A^{s}=O\) 或 \(B^{s}=O\) ．
\end{itemize}

\subsection{题目8}
8．已知 \(A B-B A=a A\) ．证明：

\begin{itemize}
  \item[（1）] 若 \(a \neq 0\) ，则存在可逆矩阵 \(P\) 使得 \(P^{-1} A P, P^{-1} B P\) 同时为上三角矩阵。
  \item[（2）] 若 \(a=0\) ，则 \(r(A+B)+r(A B) \leq r(A)+r(B)\) 。
\end{itemize}
\end{document}
